\documentclass[10pt,a4paper]{article}
\usepackage{amsmath,amssymb,amsthm,graphicx}
\usepackage{hyperref}
\usepackage{latexsym}
\usepackage{amssymb}
\usepackage{amsthm}
\usepackage{ulem}
\usepackage{enumitem}
\setlist{itemsep=0.4pt,topsep=0pt}
\parindent 0pt
\parindent 0pt
\parskip 12pt
\newcommand{\HCF}{\mathrm{hcf}}
\newcommand{\ZZ}{\mathbb{Z}}
\newcommand{\rk}{\mathrm{rk}}
\DeclareMathOperator{\len}{length}
\DeclareMathOperator{\bl}{b}
\DeclareMathOperator{\cl}{cl}
\DeclareMathOperator{\di}{d}
\DeclareMathOperator{\Cl}{Cl}
\DeclareMathOperator{\acl}{acl}
\DeclareMathOperator{\tp}{tp}
\DeclareMathOperator{\Th}{Th}
\newcommand{\Cor}{\mathcal{C}_\mathrm{or}}
\newcommand{\Cun}{\mathcal{C}_\mathrm{un}}
\newcommand{\mons}{\mathbb{M}}
\newcommand{\strong}{\mathrm{strong}}
\newcommand{\axioms}{\Delta_{\mathrm{FSTS}}}
\newtheoremstyle{note}% hnamei
{8pt}% hSpace abovei
{14pt}% hSpace belowi
{\normalfont}% hBody fonti
{}% hIndent amounti1
{\bfseries}% hTheorem head fonti
{ }% hPunctuation after theorem headi
{.5em}% hSpace after theorem headi2
{}% hTheorem head spec (can be left empty, meaning ‘normal’)i
\theoremstyle{note}
\newtheorem{thm}{Theorem}[section]
%\newtheorem{theorem}{Theorem}
\newtheorem{lemma}[thm]{Lemma} 
\newtheorem{df}[thm]{Definition}
\newtheorem{rmk}[thm]{Remark} 
\newtheorem{prop}[thm]{Proposition}
\newtheorem{fact}[thm]{Fact}
\newtheorem{cor}[thm]{Corollary}
\newtheorem{conj}[thm]{Conjecture}
\newtheorem{ex}[thm]{Example}
\newtheorem{qu}[thm]{Question}
\newtheorem*{notation}{Notation} 
\title{An elementary embedding of a free system into itself is surjective}
\author{Silvia Barbina, Enrique Casanovas}
\date{arxiv:2411.13723v3; Creative Commons Attribution 4.0 International}
\begin{document}
\maketitle
\noindent\textit{Editorial note.} The counted statement is the author's environment at source0.tex lines 1076-1077 and its proof environment opens at line 1078; the pairing is by adjacency and the retained proof range is the whole environment.\par\medskip
\noindent\textit{Editorial note.} The preamble is the author's own package lines, macro block and theorem-environment declarations (source lines 4-65); the author's optioned amsmath line and the fleqn-dependent layout line are not retained.\par\medskip
\noindent\textit{Editorial note.} The source carries no label-number dump and no counter values are imposed, so the printed numbering is the wrapper's default; identity is the author label, or the pinned source position where the author gave none.\par\medskip
\noindent\textit{Editorial note.} Metadata provenance: the source record of the pinned copy and the licence evidence file are both present; the licence is CC BY 4.0 with status supported, and the card's source\_version matches the value recorded in that evidence file.\par\medskip
\noindent\textit{Editorial scope.} One original result of the pinned author manuscript, counted once: elementary embeddings of the free system into itself are surjective (author label April\_11\_2024\_5).\par\medskip
\noindent\textit{Allowed standard prerequisites (editorial).}\begin{itemize}
\item Free system on finitely many generators. F\_n is freely generated by a set of n elements, so its generated subsystems correspond to subsets of the free base.. It is the object of the statement.
\item Elementary embedding. An embedding preserving the truth of all formulas with parameters.. It is the hypothesis.
\item Surjective endomorphism is an automorphism. Every surjective homomorphism of F\_n onto itself is an automorphism.. It converts surjectivity into the automorphism conclusion.
\end{itemize}\medskip
	\begin{rmk}\label{April_11_2024_4} Let $F_n$ be the free STS with $n\geq 3$ generators. Every surjective homomorphism $f:F_n\rightarrow F_n$  is an automorphism of $F_n$.
	\end{rmk}

	\begin{prop}\label{April_11_2024_5} Let $F_n$ be the free STS with $n\geq 3$ generators. Every elementary embedding $f:F_n\rightarrow F_n$ is surjective and, therefore, is an automorphism of $F_n$.
	\end{prop}

	\begin{proof} Let  $F_n=\langle \bar{a}\rangle$, where $\bar{a}=a_1,\ldots,a_n$ is a free base,  and let $a_i^\prime =f(a_i) = t_i(\bar{a})$  for every   $i=1,\ldots,n$, where $t_i$ is a corresponding term. Then  $F_n$ satisfies
		$$\exists x_1\ldots x_n  \bigwedge_{i=1}^n  a_i^\prime = t_i(x_1,\ldots,x_n).$$
		Since $f(F_n)\preceq F_n$, the same sentence holds in $f(F_n)$ and, by isomorphism,  $F_n$ satisfies
		$$\exists x_1\ldots x_n  \bigwedge_{i=1}^n  a_i= t_i(x_1,\ldots,x_n).$$
		Choose $\bar{b}= b_1,\ldots,b_n\in F_n$  such that   $a_i= t_i(\bar{b})$  for every $i=1,\ldots,n$. Then  $F_n=\langle b_1,\ldots,b_n\rangle$  and $b_1,\ldots,b_n$ are independent by Remark~\ref{April_11_2024_4}. Now consider the automorphism $g:F_n\rightarrow F_n$  defined by  $g(b_i)=a_i$  for  $i=1,\ldots,n$. Observe that  for every $i$
		$$g(a_i)= g(t_i(\bar{b}))= t_i(g(\bar{b}))=t_i(\bar{a})= a_i^\prime.$$
		This implies that $f=g$ and, therefore, that $f$ is an automorphism.
	\end{proof}

\end{document}
